\documentclass[11pt]{article}
\usepackage[a4paper,margin=26mm]{geometry}
\usepackage{amsmath,amssymb,mathtools}
\usepackage{unicode-math}
\setmainfont{Latin Modern Roman}
\setmathfont{Latin Modern Math}
\usepackage{xurl,hyperref}
\hypersetup{hidelinks,pdftitle={UnsolvedMath Problem 10300050}}
\newcommand{\sref}[2]{\hyperlink{source:#1:#2}{[S#2]}}
\newcommand{\cataloguescope}[1]{#1}
\begin{document}
% UnsolvedMath ID 10300050; source catalogue code AMR-102-0050
\setcounter{section}{10300049}
\section{Foliated Teichmüller theory — Question 11.4}\label{10300050}
{\small\sffamily UnsolvedMath ID 10300050 · Topology}
\subsection{Definitions and mathematical statement}

\paragraph{Shared notation.}$\mathbb N=\{1,2,\ldots\}$, and $\mathbb Z,\mathbb Q,\mathbb R,\mathbb C$ denote the integers, rationals, reals and complex numbers. Any other conventions are specified locally.


A codimension-one foliation $\mathcal F$ decomposes a $3$-manifold $M$ into surface leaves, locally as the horizontal slices of a product chart. In the terminology of the source, it is taut if a closed curve transverse to $\mathcal F$ meets every leaf. A minimal foliation means that every leaf is dense; a minimal surface means a surface with vanishing mean curvature for the specified Riemannian metric. These two uses of minimal are distinct. Equip leaves with a conformal structure. Let $(\gamma_i)_{i\in I}$ be a collection of transverse circles whose union meets every leaf and $n_i\ge0$ integers. The divisor notation $D=\sum_i n_i\gamma_i$ specifies poles of order at most $n_i$ at each leaf intersection with $\gamma_i$. Write $\mathscr O(D)$ for functions leafwise holomorphic off these circles and leafwise meromorphic with these pole bounds.
\paragraph{Statement recorded in the supplied source.}
Determine how $\mathscr O(\sum_i n_i\gamma_i)$ varies with the transverse circles $\gamma_i$. In particular determine effects of topological modifications of these circles, including crossing changes and cabling, while retaining valid transverse data. Include the function spaces themselves, not only dimensions in a particular example. \sref{10300050}{2}
\subsection{Short English statement}
Describe how spaces of leafwise holomorphic functions with prescribed transverse poles change when the transverse circles move, cross or are cabled.
\subsection{Sources}
\begin{description}
\item[\hypertarget{source:10300050:1}{S1}] ULAM.AI and the UnsolvedMath Contributors, \emph{UnsolvedMath: A Curated Collection of Open Mathematics Problems}, record 10300050 (AMR-102-0050). CC BY 4.0. \url{https://huggingface.co/datasets/ulamai/UnsolvedMath}.
\item[\hypertarget{source:10300050:2}{S2}] Danny Calegari, \emph{Problems in foliations and laminations of 3-manifolds}, version 0.78 (2002), Question 11.4; Definitions 1.1--1.10 and local definitions and remarks. \url{https://arxiv.org/abs/math/0209081v1}.
\end{description}
\label{end:10300050}
\end{document}
